%  LaTeX support: latex@mdpi.com 
%  For support, please attach all files needed for compiling as well as the log file, and specify your operating system, LaTeX version, and LaTeX editor.

%=================================================================
\documentclass[information,article,submit,pdftex,moreauthors]{Definitions/mdpi} 

%=================================================================

% MDPI internal commands
\firstpage{1} 
\makeatletter 
\setcounter{page}{\@firstpage} 
\makeatother
\pubvolume{1}
\issuenum{1}
\articlenumber{0}
\pubyear{2023}
\copyrightyear{2023}
%\externaleditor{Academic Editor: Firstname Lastname}
\datereceived{6 March 2023} 
\daterevised{11 April 2023} 
\dateaccepted{} 
\datepublished{} 
\hreflink{https://doi.org/} % If needed use \linebreak
%\doinum{}

%=================================================================
% Add packages and commands here. The following packages are loaded in our class file: fontenc, inputenc, calc, indentfirst, fancyhdr, graphicx, epstopdf, lastpage, ifthen, lineno, float, amsmath, setspace, enumitem, mathpazo, booktabs, titlesec, etoolbox, tabto, xcolor, soul, multirow, microtype, tikz, totcount, changepage, attrib, upgreek, cleveref, amsthm, hyphenat, natbib, hyperref, footmisc, url, geometry, newfloat, caption

\graphicspath{{figures/}} % path to folder with figures
\usepackage{subcaption}
\usepackage{listings}
\usepackage{xcolor}
\usepackage{gensymb} % for \textdegree
\usepackage{tabularx}
\usepackage{makecell}
\usepackage{array}
\newcolumntype{?}{!{\vrule width 1pt}}

%=================================================================
% Full title of the paper (Capitalized)
%\Title{Land Features Recognition in Qena Bend of the Nile River, Egypt Using Unsupervised Classification of the Landsat Images by R}
\Title{Recognizing the Wadi Fluvial Structure and Stream Network in Qena Bend of the Nile River, Egypt on Landsat 8-9 OLI Images}

% MDPI internal command: Title for citation in the left column
\TitleCitation{Recognizing the Wadi Fluvial Structure and Stream Network in Qena Bend of the Nile River, Egypt on Landsat 8-9 OLI Images}

% Author Orchid ID: enter ID or remove command
\newcommand{\orcidauthorA}{0000-0002-5759-1089} % Polina Add \orcidA{} behind the author's name
\newcommand{\orcidauthorB}{0000-0002-6461-1551} % Olivier Add \orcidB{} behind the author's name

% Authors, for the paper (add full first names)
\Author{Polina Lemenkova $^{1}$*\orcidA{} and Olivier Debeir $^{1}$\orcidB{}}

%\longauthorlist{yes}

% MDPI internal command: Authors, for metadata in PDF
\AuthorNames{Polina Lemenkova and Olivier Debeir}

% MDPI internal command: Authors, for citation in the left column
\AuthorCitation{Lemenkova, P.; Debeir, O.}
% If this is a Chicago style journal: Lastname, Firstname, Firstname Lastname, and Firstname Lastname.

% Affiliations / Addresses (Add [1] after \address if there is only one affiliation.)
\address{%
$^{1}$ \quad Laboratory of Image Synthesis and Analysis (LISA), École polytechnique de Bruxelles (Brussels Faculty of Engineering), Université Libre de Bruxelles (ULB). Building L, Campus du Solbosch, ULB -- LISA CP165/57, Avenue Franklin D. Roosevelt 50, Brussels 1050, Belgium.
}

% Contact information of the corresponding author
\corres{Correspondence: polina.lemenkova@ulb.be; Tel.: +32 471 86 04 59 (P.L.)}

% Abstract
\abstract{With \textcolor{blue}{methods for processing }remote sensing data  \textcolor{blue}{becoming widely available}, the ability to  \textcolor{blue}{quantify }changes in spatial data and \textcolor{blue}{to evaluate the } distribution of diverse landforms across target areas in the datasets becomes \textcolor{blue}{increasingly }important. One way to approach this problem is through satellite image processing. In this paper, we primarily focus on the methods of the unsupervised classification of the Landsat OLI/TIRS images covering the region of Qena governorate in Upper Egypt. The Qena Bend of the Nile River presents a remarkable morphological feature in Upper Egypt including a dense drainage network of wadi aquifer systems and plateau largely dissected by numerous valleys of dry rivers. To identify the fluvial structure and stream network of Wadi Qena region, this study addresses the problem of \textcolor{blue}{interpreting the relevant space-borne data using R, with an aim to visualise } the land surface structures corresponding to various land cover types. \textcolor{blue}{To this effect, high-resolution both } 2D and 3D topographic and geologic maps were used for \textcolor{blue}{the }analysis of the geomorphological setting of Qena region. The information was extracted from the space-borne data for comparative analysis of the distribution of wadi streams in  \textcolor{blue}{the }Qena Bend area  \textcolor{blue}{over }several years: 2013, 2015, 2016, 2019, 2022, and 2023. Six images were processed using \textcolor{blue}{computer vision methods made available }by R libraries. The results of the \emph{k}-means clustering of each scene retrieved from the multi-temporal images covering Qena Bend of the Nile River were  \textcolor{blue}{thus }compared to visualise changes in landforms caused by the cumulative effects from the geomorphological disasters and climate-environmental processes. The proposed  \textcolor{blue}{method, tied together through the use }of R scripts \textcolor{blue}{, }runs effectively and performs favourably in computer vision  \textcolor{blue}{tasks } aimed at geospatial image processing and analysis of remote sensing data.} 

% Keywords
\keyword{Africa; k-means clustering; image processing; remote sensing; unsupervised classification; programming; visualization; modelling; computer vision; cartography} 

\PACS{91.10.Da; 91.10.Jf; 91.10.Sp; 91.10.Xa; 96.25.Vt; 91.10.Fc; 95.40.+s; 95.75.Qr; 95.75.Rs; 42.68.Wt}
\MSC{86A30; 86-08; 86A99; 86A04}
\JEL{Y91; Q20; Q24; Q23; Q3; Q01; R11; O44; O13; Q5; Q51; Q55; N57; C6; C61}

%%%%%%%%%%%%%%%%%%%%%%%%%%%%%%%%%%%%%%%%%%
\begin{document}

\section{Introduction}

\subsection{\textcolor{blue}{Background}}
Satellite images are frequently used as a source of geoinformation for mapping land cover types and recognising land features. Example of such applications include thematic maps made using classification of satellite images; these include, for instance maps of mineral resources \cite{ABDELKAREEM2018682,1294518,rs11121450}, environmental assessment \cite{615837,Lemenkova202227}, land cover/land use classification \cite{516791,Lemenkova202229}, geologic analysis \cite{6723607,KAMEL2022104420} and urban mapping \cite{MOHAMED2019103507}. The success of the remote sensing data in cartographic processing and mapping is based on the effectiveness of the satellite imagery as a source of information for recognition of diverse land features and processes. 

Among the advantages of \hl{using }the high-resolution \hl{space-borne }data as a source of information\hl{, }one can mention the improved quality of spatial visualization \cite{ABDALLA20128,Lemenkova202230}, and the enabled access to remotely located places and areas otherwise inaccessible, such as deserts \cite{6350376}. For such places, using satellite \hl{imagery }presents a valuable source of information. Apart from spatial access, satellite images \hl{facilitate }the temporal analysis of environmental changes or disaster monitoring using \hl{the } comparative analysis of several scenes. For instance, the comparison of several satellite images covering target region in different time periods enables to perform time series analysis which can be used for detecting climate and environmental changes \hl{or }monitoring disaster events\hl{. A similar image to image comparison can also help assessing } risks of hazards, and to perform ecological surveillance of the target study area using multi-temporal data analysis \cite{9171725,SINGH2020423,Hashim}. 

Applications of remote sensing data in Earth sciences include monitoring the desertification and land degradation \cite{Badreldin,Kawy}, analysis of deforestation \cite{516791}, evaluation of salinization \cite{8949893} and soil erosion \cite{Gad2020}, hydrological modelling and quantification of the flooded areas \cite{TAHA2017157,ELSADEK2019377} and many more. Spectral patterns of the multispectral Landsat images are especially widely used in geosciences as a source of spatial information for land cover mapping and recognition of land features \cite{8486245}. To this end, the geometric, radiometric, and spectral precision of the Landsat scenes can be enhanced through a combination of the high-resolution panchromatic and multispectral channels by fusion methods \cite{4420687}.

Natural resource monitoring benefits from the use of satellite images which provide a powerful source of information for detecting the environmental processes \cite{Khafagy}. \hl{Furthermore}, identifying lithologic land surface structures and extracting extracted structural lineaments is possible using computer vision techniques \cite{Bishta,Kusky,Salman}. For instance, research on mineral exploration can be supported using the classified satellite images to discriminate lithological units on the colour composites of the images in various ratios \cite{Shalaby}. For geomorphological modelling, the Shuttle Radar Topography Mission (SRTM) Digital Elevation Model (DEM) can be overlaid with satellite images \cite{Liping}. Such data integration enables to identify the morphometric parameters of the river catchment and drainage area, and to examine spatial distribution of the agriculture fields with regards to the fluvial networks \cite{6265340}. 

Information derived from the classified satellite images enables to perform the assessment of geomorphological and geological risks for defining maps of hazards. This includes the problem of morphometric modelling using remote sensing data for identification of geologic hazards, as widely investigated in cartography, geoinformatics and Earth science communities. Hamdan and Khozyem \cite{Hamdan} applied the Advanced Spaceborne Thermal Emission and Reflection Radiometer (ASTER) data and Geographic Information System (GIS) techniques to delineate drainage networks in Wadi El-Mathula watershed in Central Eastern Desert. Various methods of geomorphometric modelling exist and can be applied to detect the spatial patterns of fluvial systems, or for hydrogeological modelling aimed at the analysis of groundwater potential \cite{Issawi,Abdalla2020,Gaber}. The geometry and morphology of the relief, the characteristics of topographic parameters such as slope, aspect and hillshade can be extracted as morphometric characteristics for analysis of watershed using 3D mapping \cite{Lemenkova202301}. 

Extracting the information from the remotely sensed data can be used for complex hydrological analysis, for instance, predicting wadi flash floods using morphometric parameters that represent the topographic and drainage characteristics of the basins using the GIS \cite{w9070553}. Further, such information can be used for the analysis of wadi networks to define basin boundaries and drainage networks, to evaluate stream channel density, and to model slope and flow directions. Moreover, it is useful for retrospective modelling of the hydrological processes in the past to better analyse geologic setting in \hl{the present, } using robust remote sensing data. Such approach enables deeper insights in the structure of the land surface visible from the \hl{space-borne } data \cite{abdelkareem2013integration,AbdelkareemFarouk}. 

\hl{Also}, the essential geographic information collected from the satellite imagery is useful for operative flash flood monitoring\hl{, through integrating } data on physiographic features. Other applications of the remote sensing data include modelling lithological, geological, and structural features \cite{BESHR2021999}, mapping mineral resources \cite{Kamal} and detecting wadi channels based on the combination of the topographic maps and satellite images \cite{Moubark}. 

\subsection{\textcolor{blue}{Motivation and Objectives}}

\hl{In view of the advantages of the satellite images for Earth and environmental studies, a question of effective methods of processing these data is arising. A variety of existing works have explored the use of GIS for remote sensing data analysis. The use of ENVI GIS} \cite{ABDELSADEK2022815,GHEBREZGABHER201637,ABDELMALIK2020263} \hl{well illustrates the traditional methods of remote sensing data processing by the use of supervised classification. More examples of the existing studies have used the combination of software for satellite image processing. A prominent example of such an approach is presented by} \cite{KAMEL2016343} \hl{wherein the use of several GIS is assigned to the relevant tasks in image processing: ArcGIS -- for interpretation of the geological lineaments, Erdas Imaging -- for data preprocessing, ILWIS GIS and ENVI GIS -- for creating band colour composites and information extraction. Moreover, the integration of the ArcGIS and ENVI GIS is presented in} \cite{ELDOSOUKY20201008,GABR202211} \hl{for processing aeromagnetic and geologic data for mapping the structural complexity and mineralization in Southern Eastern Deserts of the Upper Egypt.}

\hl{Other studies have used Erdas Imagine for processing multi-temporal remote sensing data, e.g.,} \cite{BAKR2010592,NASSERMOHAMEDEID202066} \hl{monitored land cover changes and computed vegetation indices. Similarly,} \cite{ALIBIK2022113} \hl{presented the application of the Erdas Imaging for geological analysis through identifying rocks by their spectral signatures.} \cite{ALIBIK2022593} \hl{applied ArcGIS for detecting lineaments on DEM and Sentinel-2 images for geologic analysis of southern Sinai;} \cite{HOSSEN2016951} \hl{combined ArcGIS and Erdas Imagine for multi-temporal image processing aimed at environmental impact assessment. Such methods correspond to the aim our approach in terms of using satellite image processing for monitoring landscapes of Egypt, however, they are largely tied to specific niches and problem formulation such as land cover changes, stratigraphic and petrographic analysis and identifying rocks by the use of remote sensing data. Nevertheless, while these algorithms benefit from the use of the conventional interface in the traditional GIS software, they are limited by the largely restricted functionality of all these programs. }

\hl{In contrast to these and similar existing studies, here }we propose an R-based \hl{use case of }satellite image processing to address the problem of \hl{using computer vision } for extracting geoinformation from the remotely sensed \hl{space-borne }data. The main motivation of this work is to demonstrate that land surface types can be extracted from a short time series of multispectral Landsat 8-9 OLI/TIRS images by k-means clustering \hl{, using libraries of the }R language. This work reports an empirical comparison of the classification of the six Landsat scenes for years 2013, 2015, 2016, 2019, 2022 and 2023 to detect changes in the identified land surface types for target study region on Qena Bend, Upper Egypt. Furthermore, the results report a comparison of the detected classes by Landsat bands in false colour composite (5-4-3 channels corresponding to the Near-Infrared, Red and Green Bands in the Landsat OLI/TIRS images). The results on the performed classification support the argument that satellite images can be integrated out to achieve a recognition of the land surface types in a complex drainage network of Wadi Qena, Upper Egypt, and this paper demonstrates \hl{a }practical approach to achieve satellite image processing \hl{using scripts written in the }R programming language.

%\pagebreak
\section{Prior Work}

\hl{The }Qena Bend is one of the most notable morphological features of the Nile River in Upper Egypt\hl{, with }a special characteristic curvature in its geometry \hl{as shown in }Figure \ref{fig01}. The structural geomorphology of the Qena Bend largely reflects the general geologic setting of the Nile River and major processes of its evolution \cite{Said1981,said2017geology}. The Valley of the Nile River has been cut during late Miocene time of the Neogene period and formed later on as a results of the regional and local tectonics. The Qena Formation includes the gravel deposits with limestone cobble and heavily weathered soil typical for old terrace deposits \cite{paulissen1987egyptian}. The remaining sediments, outcrops and paleospecies in fossils are identified from the Mesozoic Era (Cretaceous period) \cite{ABDELHAMID2014138,SOLIMAN1986111} and Paleogene (Paleocene and Eocene epoches) \cite{MANDUR2022104594,sehim1993cretaceous}. They are reflected in the Cenozoic sediments accumulated in the lacustrine-alluvial fans around the Qena Bend due to the Nile River evolution in Neogene, Figure \ref{fig02}.

\begin{figure}[H]
%\centering
	\hspace{-40pt}\includegraphics[width=16.5 cm]{Fig_01.jpg}
	\caption{Topographic map of Egypt. Target study area of \hl{the }Qena Bend region is shown in red rotated square. Software: GMT v. 6.1.1. Data source: GEBCO/SRTM. Cartography source: authors.
	\label{fig01}}
\end{figure}

Apart from the geology and hydrology of the Nile River, the region of Qena Bend is influenced by the presence of the Eastern Desert which has a special climate. Its special feature includes the occasional extreme flash floods -- one of the most severe natural disasters. Flash floods occur annually, mostly during spring and autumn, and are caused by the torrential rainfalls followed by the disastrous surface runoff thereafter \cite{Moawad}. Such climate processes, aggravated by an interrelated complex of fluvial networks in Wadi Qena catchment area, as briefly explained above, presents favourable conditions for disastrous flash floods. During such events, a complex of the ephemeral rivers dry during most of the year is being filled with water from heavy showers and rains \cite{ElFakharany}.

The cumulative effects from \hl{the }wadi network and arid and semi-arid climate of the Eastern Desert region \hl{result }in occasional flash flood events and irregular sand dust storms with different durations and intensities. Flash floods \hl{occur }in diverse regions across  Egypt including Upper Egypt with Eastern Desert \cite{ElHaddad,Elsadek}, southern Red Sea Coast \cite{Abdalla2014}, Sinai Peninsula \cite{ELOSTA2021101522} and Gulf of Suez \cite{AbdelLattif}. The social consequences of flash floods include damaged infrastructure, transport ways and buildings, and occasional victims \cite{Moawad}. El-Magd et al. report the catastrophic events of flash floods in Qena governorate recorded in three consecutive events in 2014--2016 \cite{ElMagd}. 

\begin{figure}[H]
%\centering
	\hspace{-60pt}\includegraphics[width=16.5 cm]{Fig_02}
	\caption{Geologic units of Egypt and the Nile River. Software: QGIS v. 3.2.2. Data source: USGS. Cartography source: authors.
	\label{fig02}}
\end{figure}

As a result of the geologic evolution, the geomorphology of the Qena Bend region present a plateau largely dissected by numerous valleys of wadi and geologic lineaments which dissect the cliffs of the Nile River and basement plateau composed of the limestone and chalks \cite{BANDEL1987427}. This result in a specific geomorphic pattern of the Qena Bend surroundings which is notable for a complex drainage network of wadi aquifer systems \cite{Alkholy,AbdelMoneim,HUSSIEN201773}. Such drainage network presents a typical feature of the Eastern Desert which consist of numerous fluvial basins that drain the rainwater towards the Nile River or the Red Sea \cite{Abdel}.

Major geomorphic structure of the Nile valley follows the faults oriented parallel to the Red Sea along the central current of the Nile and the Gulf of Suez in the north \cite{Ali2022,Youssef2003StructuralSO} with the relics of the pre-late Miocene streams remaining in fluviatile sand and gravel sediments of the Nile delta in Lower Egypt \cite{Said1981a}. The discovery of paleorivers in the Sahara by radar imaging systems \hl{revealed } the orientation of the paleochannels and their flow directions \cite{ABDELKAREEM201235}. During the Neogene, it is argued \cite{PHILOBBOS2015194} that the ancient Qena Lake broke up the Nile River course which might have affected its present special form and specific geometric curvature, Figure \ref{fig03}.

The occurrence of flood hazards in Upper Egypt point at the significance of the application of geoinformation for visualization and mapping regions prone to flood risk. The remote sensing data and advanced methods of spatial data processing present robust approaches to evaluation of relief and extracting information on relief and topographic characteristics from the satellite images using methods of computer vision. Such approaches can be used for \hl{instance, for }vulnerability mapping or identification of areas with high geomorphological hazards and \hl{exposure }to floods \cite{nhess-9-751-2009}.

\begin{figure}[H]
\centering
	\includegraphics[width=14.5 cm]{Fig_03}
	\caption{Enlarged fragment of the 3D perspective view of Qena Bend, Nile River, Upper Egypt. Software: GMT v. 6.1.1. Data source: GEBCO/SRTM. Cartography source: authors.
	\label{fig03}}
\end{figure}

%%%%%%%%%%%%%%%%%%%%%%%%%%%%%%%%%%%%%%%%%%
\section{Materials and Methods}

\subsection{Data}

Multispectral satellite images such as Landsat TM/ETM+, and OLI/TIRS) provide a valuable source of remote sensing data widely used in environmental studies of Egypt \cite{Elfadaly,Sultanetal1986,Sultanetal1987}. The reason for \hl{selecting } the Landsat OLI/TIRS images is due to their open source availability, regular global coverage and data robustness. The Landsat 8-9 OLI/TIRS scenes have 30-m resolution in multispectral channels along a 185 km swath for each image (Bands 1 to 7, used for colour composites), systematically presented, geometric, radiometrically calibrated, georeferenced using the World Geodetic System 1984 (WGS84) datum, UTM Coordinate system (Zone 36 for this case) and terrain corrected by the source provider. Such characteristics of the Landsat 8-9 images makes them a reliable and robust source of geo-information. 

Six Landsat OLI/TIRS scenes covering the study area with minimal cloud coverage were selected for analysis and interpretation of the land cover types. \hl{The selection of different image time intervals is explained by the availability of the Landsat 8-9 OLI/TIRS with cloud-free scenes covering target area. In particular, Landsat OLI/TIRS sensor is improved against the earlier versions of Landsat products in terms of technical and spectral characteristics such as narrower spectral bands, refined radiometric resolution, better noise detection and calibration parameters} \cite{IRONS201211}. \hl{Since spectral characteristics of the older Landsat products (Landsat TM and Landsat MSS) has worse parameters, we only selected the data from the Landsat OLI/TIRS with the first images are available starting from 2013 when the first OLI/TIRS sensor was launched.} 

\hl{Given the advantages of the Landsat 8-9 OLI/TIRS sensor which provide moderate-resolution scenes with multi-spectral bands available for the dates between 2013 and 2022, they were selected for this research. The time span was selected based on the availability of the data in the Qena area that demonstrate low cloudiness in winter period with the interval of 1-3 years. The images were selected in winter period to avoid impacts from the arid climate on vegetation parameters. Thus, the images} were acquired in November for the Landsat-8 scenes for years 2013 to 2022 and January for Landsat-9 image for 2023. Technical characteristics and metadata of the Landsat satellite images used for image processing are summarised in Table \ref{tab1}. 

\begin{table}[H]
\small
\caption{Satellite images used for computing the VIs: Landsat-8 OLI/TIRS collected from the USGS. \textsuperscript{1}.\label{tab1}}
	\begin{adjustwidth}{-\extralength}{0cm}
%		\newcolumntype{C}{>{\centering\arraybackslash}X}
		\begin{tabularx}{\fulllength}{ p{1.6cm} | p{1.8cm} | p{7.0cm} | p{4.0cm} | p{2.0cm} }
			\toprule
			\textbf{Date} & \textbf{Spacecraft} & \textbf{Landsat Product ID} & \textbf{Scene ID} & \textbf{Cloudiness} \\
			\midrule
			2013/11/16 & Landsat 8 & \makecell{LC08\_L2SP\_175042\_20131116\_20200912\_02\_T1} & LC81750422013320LGN01 & 0.89 \\
			2015/11/22 & Landsat 8 & \makecell{LC08\_L2SP\_175042\_20151122\_20200908\_02\_T1} & LC81750422015326LGN01 & 0.05 \\
			2016/11/08 & Landsat 8 & \makecell{LC08\_L2SP\_175042\_20161108\_20200905\_02\_T1} & LC81750422016313LGN01 & 0.31 \\
			2019/11/17 & Landsat 8 & \makecell{LC08\_L2SP\_175042\_20191117\_20200825\_02\_T1} & LC81750422019321LGN00 & 1.16 \\
			2022/11/09 & Landsat 8 & \makecell{LC08\_L2SP\_175042\_20221109\_20221121\_02\_T1} & LC81750422022313LGN00 & 0.02 \\
			2023/01/20 & Landsat 9 & \makecell{LC09\_L2SP\_175042\_20230120\_20230122\_02\_T1} & LC91750422023020LGN01 & 0.32 \\	
			\bottomrule
		\end{tabularx}
	\end{adjustwidth}
	\noindent{\footnotesize{\textsuperscript{1} The Sensor ID is common for all the scenes: Landsat 8-9 OLI/TIRS (Operational Land Imager and Thermal Infrared Sensor), Collection 2 Level-2. Path/row parameters common for all the images: 175/42. Coverage: Qena Bend, the Nile River, Upper Egypt. Image source: the USGS}}
\end{table}

Besides the Landsat OLI/TIRS images, we also use the topographic and geologic grids for mapping the study area and analysis of the geomorphology in 2D and 3D formats. The interpretation of these ancillary data helps to analyse and verify the impact of the geologic processes and lithological structure in the temporal variations in land cover and surface structure. In such a way, this work is based on the use of the two major types of data: six satellite Landsat 8-9 OLI/TIRS images and topographic-geological datasets Figure \ref{fig04}. 

\begin{figure}[H]
\centering
	\includegraphics[width=14.0 cm]{Fig_04}
	\caption{Data sources: summary of the materials used in this study. Software: R version 4.2.2, 'Mermaid' library. Flowchart source: authors.
	\label{fig04}}
\end{figure}

Various land surface objects have different spectral reflectance corresponding to Digital Numbers (DN) of pixels within a given wavelength interval which corresponds to different Bands (or channels) of the Landsat images. Such \hl{an }important feature of the objects visible on the images enables \hl{their identification }by computer vision. This is \hl{made }possible using the recognition of their characteristics through the classification of Landsat Bands taken as RGB triplets of band combinations. In this study, we used different colour composites for visual analysis of the scenes and false colour composites for k-means clustering. The specific geometric curvature of the Nile River in \hl{the }Qena Bend area is perfectly visible in the Landsat images in natural colours, Figure \ref{fig05}. Here it is possible to discriminate the floodplain of the Nile River in bright green colour and a dark blue thread of the river itself contrasting against the surroundings of sands of the Western and Eastern Deserts. 

The areas of dark crimson represent massifs of bare soil\hl{, while the }ivory colour indicate a complex wadi network with its typical dendritic pattern contouring the valleys, as visible in Figure \ref{fig05}. High-resolution topographic data were used for inspection of the relief in the region and visualising the geomorphology in 2D and 3D modes, while the remote sensing data were used for comparative analysis of changes in the wadi channels in the Qena Bend area and surroundings\hl{, both }by image processing and k-means clustering algorithms using R \cite{RCoreTeam}. 

\begin{figure}[H]
\makebox[1.00\linewidth][r]{%
	\begin{subfigure}[b]{.40\textwidth}
		\centering
			\includegraphics[width=0.87\textwidth]{Fig_05a.jpg}
		\caption{2013}
	\end{subfigure}%
	\begin{subfigure}[b]{.40\textwidth}
		\centering
		\includegraphics[width=0.87\textwidth]{Fig_05b.jpg}
		\caption{2015}
	\end{subfigure}%
	\begin{subfigure}[b]{.40\textwidth}
		\centering
		\includegraphics[width=0.87\textwidth]{Fig_05c.jpg}
		\caption{2018}
	\end{subfigure}%
}\\
\vfill \vspace{1mm}
\makebox[1.00\linewidth][r]{%
	\begin{subfigure}[b]{.40\textwidth}
		\centering
			\includegraphics[width=0.87\textwidth]{Fig_05d.jpg}
		\caption{2020}
	\end{subfigure}%
	\begin{subfigure}[b]{.40\textwidth}
		\centering
		\includegraphics[width=0.87\textwidth]{Fig_05e.jpg}
		\caption{2021}
	\end{subfigure}%
	\begin{subfigure}[b]{.40\textwidth}
		\centering
		\includegraphics[width=0.87\textwidth]{Fig_05f.jpg}
		\caption{2022}
	\end{subfigure}%
}\caption{Qena Bend of the Nile River, Egypt visible on Landsat 8-9 OLI/TIRS images in natural RGB colour for six years: (\textbf{a}) 2013/12/20, (\textbf{b}) 2015/12/26, (\textbf{c}) 2018/02/21, (\textbf{d}) 2020/01/22, (\textbf{e}) 2021/12/18, (\textbf{f}) 2022/12/29.}\label{fig05}
\end{figure}

%----------- subsection ----------->
\subsection{Methods}

\subsubsection{Research concept and advantages}

\hl{The concept of the performed research include data selection and collection from the USGS repository, data processing by clustering method, analysis and interpretation of the obtained results. After data capture, cartographic visualization and image preprocessing, clustering is the major step in the research workflow. The essential approach of clustering is that it partitions the dataset into several clusters (or groups) of pixels using algorithms of data partition. Diverse types of clustering techniques exist in data analysis own strengths and weaknesses, due to the complexity of information} \cite{XuTian}. \hl{The examples of clustering include the algorithms that identify centroids, density or distribution in a dataset which enable to split the data into several groups according to the distance of each particular pixel to the centre of the cluster. Currently, clustering techniques are widely used in programming tools including machine learning} \cite{CHANDER2023179}.

\hl{In spatial data analysis, clustering is used for image analysis as a tool for unsupervised classification. Clustering algorithm separates pixels on the image into several groups (or clusters) according to their spectral signatures and assigns those pixels to the defined clusters. The most well-know algorithms of clustering is the K-Means. Other examples include Hierarchical Clustering} \cite{ZHIYING2023101935,XIE2023109230}, \hl{partition clustering} \cite{ZEIN2023100731}, \hl{fuzzy clustering} \cite{ROY2012452}, \hl{mean shift} \cite{GUO201828,BECK2019128}, \hl{density-based clustering} \cite{CHENG2023170,MAHESHWARI2023109341}, \hl{Model-based clustering} \cite{YOU2022107457}, \hl{Density-Based Spatial Clustering of Applications with Noise (DBSCAN)} \cite{OUYANG2022688,PRANATA2023319,WU202257}. 

\hl{In this study, we used the K-means algorithm of R which defines clusters of pixels on the image based on their similar features and is embedded in the R programming language} \cite{FAVERO2023193}. \hl{The concept of K-means clustering includes the iterative process of evaluation of pixels' values wherein each point is being assigned to the specific cluster group corresponding to the defined number of clusters. This approach is straightforward, implemented by the 'raster' package and accurately classifies the data on the satellite image. K-means performs the simplification of the large massifs of pixels on the image through partitioning them into groups (e.g., land cover types). The advantages of the K-means concept include easily adaptability to new datasets, and relatively straightforward algorithm of implementation in R which enables to process the dataset in a few seconds.}

\subsubsection{Data processing workflow}

The 2D and 3D maps showing the topography and geomorphic structure of the Wadi Qena region were prepared using the Generic Mapping Tools \cite{Wessel} by the existing methods of cartographic scripting \cite{Lemenkova202228}. The general workflow in R is summarised in Figure \ref{fig06} with listed libraries, and includes the following steps of data processing. 

\begin{figure}[H]
\centering
	\includegraphics[width=12.0 cm]{Fig_06}
	\caption{Methodological workflow of the processes and research steps applied in this study. Software: R version 4.2.2, 'DiagrammeR' library. Flowchart source: authors.
	\label{fig06}}
\end{figure}

Specifically, the workflow included the following steps (the examples below are given for the image 2013 and repeated for all the images, respectively). First, the images were read in to the RStudio environment as a stack of \hl{TIFF images (} .tif files\hl{) }using the R command \emph{Landsat2013 <- list.files()}. A \hl{subset }list of the file names was generated within a working folder on the computer. Then, the 'SpatRaster' object which reads the raster data spatial structure was created as follows: \emph{landsat <- rast(Landsat2013)}. The band designations for the Landsat RGB triplets were performed using command \emph{landsatRGB <- landsat[[c(5,4,3)]]} (here, the case of false colour composite).  

Afterwards, the image was visualised using the command \emph{plotRGB(landsatRGB, r=1, g=2, b=3, axes=FALSE, stretch="lin")} of the 'raster' library. The classification was based on the k-means clustering algorithm. Running the classification was done using the \hl{following command}: \emph{unC <- unsuperClass(landsatRGB, nSamples = 100, nClasses = 10, nStarts = 5)}. Here the number of classes was defined as 10, \hl{for a good separation of }the land cover classes in the Wadi Bend region. We also tried higher number of classes, however, they resulted in more coarse results. For instance, 5 and 7 classes merged several land cover classes that were evidently different, while 20 classes, in contrast, presented unnecessary partition of the image into subclasses that should be jointed. The colour palette was defined using 'RColorBrewer' library by \emph{colors <- jet(10)}, and the visualization of the image was done using the command: \emph{plot()}. The legend was then added using the command \emph{legend()}.

The triplets of Bands 5 (Near Infrared), 4 (Red) and 3 (Green) were selected for the following reasons. Healthy vegetation has a high absorption in Red (Band 4) and high reflectance in Green (Band 3) and NIR (Band 5). Therefore, such a combination helps to better discriminate vegetation areas in the floodplain of the Nile from soil, sandy and desert areas in the Western and eastern Deserts as well as Wadi valleys. In contrast, soil has a lesser reflectance spectra in NIR because its reflectance changes along with the wavelength\hl{, which is why }such as combination of bands was used \hl{instead for the calculation and clustering }of the pixels of various land cover classes based on their distinct spectral signatures. \hl{Additionally, all }the maps derived from the clustering of the satellite images were georeferenced automatically with metadata identified by R library 'raster': WGS84 datum, UTM coordinate system Zone 36.

Nevertheless, \hl{the }spectral signature of land surface differs according to the soil and rock types. Geochemical properties of soil are reflected in the presence of organic carbon, moisture percentage, level of salinity and relevant environmental characteristics of agricultural lands (contents of sulphides, fertilisers, pesticides, etc.). Likewise, reflectance spectra of diverse rock types \hl{differ }significantly, which enables to discriminate them on the satellite images using algorithms of computer vision. Therefore, we presented various combinations of the Landsat OLI/TIRS bands for analysis of land surface types on the images, as shown in Figures \ref{fig07}, \ref{fig08} and \ref{fig09}.

%%%%%%%%%%%%%%%%%%%%%%%%%%%%%%%%%%%%%%%%%%
\section{Results}

In this paper, the classified land cover maps were plotted by unsupervised classification of k-means clustering methods and the statistics on pixel assignments is summarised in Table \ref{tab2}. \hl{In a discrimination of land cover types, spectral signature of the satellite images are taken as markers for distinguishing the landscape features based on their texture, shape and size. In this case, colour composites created based on band combination of the Landsat images are included in the data processing and clustering. }The results of the clustering include the defined 10 classes best representing the land surface structure and corresponding to the following items modified and generalised from the existing study \cite{Kamel}: \emph{1) Floodplain of the Nile River; 2) Silt sediments; 3) Wadi deposits; 4) Gritty and gravel soil;  5) Sandy soils; 6) Fine sand, silty and clay soil; 7) Gravel and stony soil; 8) Terrace soils and stony debris; 9) Limestone foothills; 
10) Limestone rock land}. 

\hl{The information was retrieved from the analysis of clustering implemented on the Landsat satellite images that represents the multi-component structure of the Earth's surface and different types of the landscapes. }The Landsat 8-9 OLI/TIRS images were differentiated \hl{using clustering techniques which enabled to analyse the images in terms of spatial distribution of pixels with different values according to their spectral signatures. The }land cover types of the Qena Band region were recognised automatically by R for the years 2013, 2015, 2016, 2019, 2022, and 2023. The dimensions of all the images are identical for all the Landsat scenes and correspond to the 7801 rows, 7651 columns, and 59,685,451 cells (or pixels) with resolution of 30x30 m, i.e., each pixel on the image corresponds to the 30 m patch size on the land surface. Thus, the number of pixels computed for each ID class gives the area of the land cover types with respective changes over years. The statistics on pixels assigned to the land cover ID classes is based on the spectral signature of the DN of the pixels, Table \ref{tab2}. \hl{Correspondingly}, Tables \ref{tab3}, \ref{tab4} and \ref{tab5} show the detected numbers of pixels for each land cover ID class  in Qena Band region, for quantitative pairwise comparison for the years 2013 with 2015, 2016 with 2019, and 2022 with 2023, respectively. 

\begin{table}[H]
\small
\caption{Results of the k-means classification of the Landsat-8 OLI images for Qena Bend area\textsuperscript{1}.\label{tab2}}
	\begin{adjustwidth}{-\extralength}{0cm}
		\begin{tabularx}{\fulllength}{ C | C | C | C | C | C | C }
			\toprule
			\textbf{Class ID} & \textbf{2013} & \textbf{2015} & \textbf{2016} & \textbf{2019} & \textbf{2022} & \textbf{2023} \\
			\midrule
			1 & 6237758.6 & 25058427 & 11830859 & 11050520 & 33773365.2 & 22414794 \\
			2 & 26530972.6 & 3829403 & 6847520 & 7687698 & 21467288.5 & 4080038 \\
			3 & 46356528.0 & 23160871 & 13056437 & 31925587 & 466472.5 & 13918989 \\
			4 & 10314256.8 & 7458695 & 12594525 & 3359878 & 3547592.0 & 1863680 \\
			5 & 7861914.8 & 9909609 & 3012708 & 11101016 & 7714730.1 & 6083803 \\
			6 & 625035.3 & 4568955 & 2195275 & 14290386 & 9346294.4 & 14743313 \\
			7 & 22694195.3 & 3019950 & 12277811 & 14898378 & 6978872.8 & 6370890 \\
			8 & 5804917.1 & 1942973 & 10144633 & 5386949 & 16906933.2 & 6053497 \\
			9 & 7320502.4 & 13450745 & 3220846 & 5501193 & 8429786.2 & 11952344 \\
			10 & 6121676.5 & 10228437 & 4195904 & 10594134 & 11161798.4 & 10517903 \\	
			\bottomrule
		\end{tabularx}
	\end{adjustwidth}
	\noindent{\footnotesize{\textsuperscript{1} Sum of squares by clusters within computed cluster centroids.}}
\end{table}

\hl{The land cover class 'floodplain' decreased on 5.5\% from 2013 (ID class 2 with 26530972.6 pixels) to 2015 (ID class 1 with value 25058427), then slightly increased in the end of the study period from 21467288.5 to 22414794 of the assigned pixels, i.e., on 4.22\%, Table} \ref{tab2}. \hl{The land class 2 'silt sediments'  showed changes from 2016 with 10144633 assigned pixels against 10594134 pixels in 2019, which demonstrates an increase to 4.24 \%. The land class 3 'wadi deposits' demonstrated the following changes in the distribution of pixels from 6237758.6 in 2013 to 6847520 in 2015, that is, the increase on 8.9\%. From 2022 to 2023, the values changes from 6978872.8 to 6053497, i.e., 13.2\%. A slight increase in the same land cover class was further noted in 2022 with 11161798.4 assigned pixels, followed by a slight decrease to 10517903 in 2023, which resulted in decline in 5.77\%.}

\hl{A slight increase in 8.73\% in the land cover class 8 'terrace soils and stony debris' was noted from 2015 with 3829403 pixels to 4195904 pixels in 2016. The land cover classes 'limestone foothills' remained with comparable values in 2015 with 13450745 pixels against 12277811 in 2016 which is a decrease in 8.72 \%, and 11101016 in 2019, i.e. on 9.58\%. Finally, the land cover class 10 'limestone rock land' experienced recent changes from 2022 with 8429786.2 values to 6083803 in 2023 which might be caused by the soil erosion processes.}\hl{To compare the land cover types over Qena, we can start with Figure} \ref{fig07} \hl{and analyse changes in land cover classes and how the pixel classes are assigned as summarised in Table} \ref{tab3}. 

Figure \ref{fig07} shows several band combinations for colour composites. \hl{Notably}, Band 4-3-2 represent a natural colour composite; Band 5-4-3 is a composition that includes an infrared (Band 5) and is therefore adjusted for identification of vegetation in the river floodplain. Band 5-7-1 contains NIR (Band 5), SWIR2 (Band 7), sensitive to radiation emission, and coastal aerosol (Band 1), sensitive to small particles, haze and also burnt areas \cite{8611806}. The area of the Nile floodplain with riparian vegetation has a strong absorption both in visible and NIR Bands, due to the pigmentation of chlorophyll which absorbs highly at wavelength in visible bands. Other factors are the physiological structure of plant canopy and moisture of leaves. Therefore, the floodplain area has a strong reflectance in all bands due to the presence of vegetation with high water content, as can be seen in Figures \ref{fig07}, \ref{fig08} and \ref{fig09}.

\begin{figure}[H]
\makebox[1.00\linewidth][r]{%
	\begin{subfigure}[b]{.40\textwidth}
		\centering
			\includegraphics[width=0.87\textwidth]{Fig_07a}
		\caption{2013: Bands \hl{4-3-2}}
	\end{subfigure}%
	\begin{subfigure}[b]{.40\textwidth}
		\centering
		\includegraphics[width=0.87\textwidth]{Fig_07b}
		\caption{2013: Bands \hl{5-4-3}}
	\end{subfigure}%
	\begin{subfigure}[b]{.40\textwidth}
		\centering
		\includegraphics[width=1.00\textwidth]{Fig_07c}
		\caption{2013: Clustering}
	\end{subfigure}%
}\\
\vfill \vspace{1mm}
\makebox[1.00\linewidth][r]{%
	\begin{subfigure}[b]{.40\textwidth}
		\centering
			\includegraphics[width=0.87\textwidth]{Fig_07d}
		\caption{2015: Bands \hl{5-7-1}}
	\end{subfigure}%
	\begin{subfigure}[b]{.40\textwidth}
		\centering
		\includegraphics[width=0.87\textwidth]{Fig_07e}
		\caption{2015: Bands \hl{6-3-2}}
	\end{subfigure}%
	\begin{subfigure}[b]{.40\textwidth}
		\centering
		\includegraphics[width=0.92\textwidth]{Fig_07f}
		\caption{2015: Clustering}
	\end{subfigure}%
}\caption{Qena Bend of the Nile River: false colour composites and clustering of the Landsat 8-9 OLI/TIRS images: (\textbf{a}) 2013 RGB in \hl{4-3-2} Bands, (\textbf{b}) 2013 RGB in \hl{5-4-3} Bands, (\textbf{c}) 2013 unsupervised classification, (\textbf{d}) 2015 RGB in \hl{5-7-1} Bands, (\textbf{e}) 2015 RGB in \hl{6-3-2} Bands, (\textbf{f}) 2015 RGB: Clustering.}\label{fig07}
\end{figure}

\begin{table}[H]
\small
\caption{Pairwise comparison of the k-means classification for Bands 5-4-3 of the Landsat-8 OLI/TIRS images for 2013 and 2015 in Qena Bend region. \textsuperscript{1}.\label{tab3}}
	\begin{adjustwidth}{-\extralength}{0cm}
		\begin{tabularx}{\fulllength}{||C?C|C|C?C|C|C||}
			\toprule
			\multirow{2}{*}{$Class ID$} & \multicolumn{3}{c}{\textbf{2013}} & \multicolumn{3}{c}{\textbf{2015}} \\ 
			\textbf{} & \emph{Band 5} & \emph{Band 4} & \emph{Band 3} & \emph{Band 5} & \emph{Band 4} & \emph{Band 3} \\
			\midrule
			1 & 21019.08 & 18348.92 & 15722.58 & 19819.50 & 17226.143 & 14765.14 \\
			2 & 21956.38 & 19148.12 & 16294.50 & 12087.00 & 10917.333 & 10115.00 \\
			3 & 26892.38 & 23105.00 & 19120.92 & 18736.17 & 9956.333 & 10142.17 \\
			4 & 14308.50 & 13163.75 & 11730.75 & 24835.46 & 21496.385 & 17990.38 \\
			5 & 19246.50 & 16942.69 & 14626.81 & 23315.56 & 20178.688 & 17043.31 \\
			6 & 17656.18 & 15405.09 & 13460.00 & 26302.78 & 22807.222 & 19156.67 \\
			7 & 23465.45 & 20280.18 & 17179.82 & 15028.50 & 13507.500 & 12109.33 \\
			8 & 16549.67 & 10211.17 & 10132.83 & 29412.00 & 25219.500 & 20731.50 \\
			9 & 22484.00 & 19481.25 & 16802.50 & 17737.46 & 15833.385 & 13783.23 \\
			10 & 24895.53 & 21740.20 & 18326.47 & 17737.46 & 18747.389 & 15990.33 \\	
			\bottomrule
		\end{tabularx}
	\end{adjustwidth}
	\noindent{\footnotesize{\textsuperscript{1} Clusters are computed for each band in the composite triplets: Band 5 (NIR), Band 4 (Red) and Band 3 (Green)}}
\end{table}

Figure \ref{fig08} demonstrates the results of the clustering for pairwise comparison of years 2016 and 2019 and the colours composites \hl{using the } following Band combinations: 5-6-4; 7-5-4, 7-6-4 and 7-5-3. False colour composites 7-5-4 and 7-6-4 (urban) include the Shortwave Infrared \hl{Band } (SWIR-2) which is optimal for \hl{the } visualization of urban areas due to a better contrast between the built-up, water areas and vegetation areas. \hl{And while composites 7-5-4 and 7-5-3 } resemble each other due to the presence of SWIR1 and Red bands, the Green (Band 3) in the \hl{second composite enables a better discrimination of } the limestone rock land \hl{using a bright blue colour, whereas in composite 7-5-4, the colour used for representing limestone rock land is very similar to that used for } the gravel and stony soil\hl{, namely } orchid and magenta respectively.

\hl{Furthermore, the area covered by the land cover class 4 'gritty and gravel soil' remained stable with only a very slight decrease in the area from 2015 with 3019950 assigned pixels to 2016 with 3012708 pixels, which is below 1\%. Moreover, the land cover class 5 'sandy soils' demonstrated an increase in the recent years from 11161798.4 to 11952344, i.e. an increase to 6.6 \%. The land cover class 6 'fine sand, silty and clay soil' has changes with a slight increase from 3220846 in 2016 to 3359878 in 2019 which is a gain 4.14\% that can be cause by the environmental effects related to the hydrology of the Nile. The land cover class 7 'gravel and stony soil' showed a slight decrease from 7687698 in 2019 to 7714730.1 in 2022, which is below 1\%. }

\begin{figure}[H]
\makebox[1.00\linewidth][r]{%
	\begin{subfigure}[b]{.40\textwidth}
		\centering
			\includegraphics[width=0.87\textwidth]{Fig_08a}
		\caption{2016: Bands \hl{5-6-4}}
	\end{subfigure}%
	\begin{subfigure}[b]{.40\textwidth}
		\centering
		\includegraphics[width=0.87\textwidth]{Fig_08b}
		\caption{2016: Bands \hl{7-5-4}}
	\end{subfigure}%
	\begin{subfigure}[b]{.40\textwidth}
		\centering
		\includegraphics[width=0.90\textwidth]{Fig_08c}
		\caption{2016: Clustering}
	\end{subfigure}%
}\\
\vfill \vspace{1mm}
\makebox[1.00\linewidth][r]{%
	\begin{subfigure}[b]{.40\textwidth}
		\centering
			\includegraphics[width=0.87\textwidth]{Fig_08d}
		\caption{2019: Bands \hl{7-6-4}}
	\end{subfigure}%
	\begin{subfigure}[b]{.40\textwidth}
		\centering
		\includegraphics[width=0.87\textwidth]{Fig_08e}
		\caption{2019: Bands \hl{7-5-3}}
	\end{subfigure}%
	\begin{subfigure}[b]{.40\textwidth}
		\centering
		\includegraphics[width=0.92\textwidth]{Fig_08f}
		\caption{2019: Clustering}
	\end{subfigure}%
}\caption{Qena Bend of the Nile River: false colour composites and clustering of the Landsat 8-9 OLI/TIRS images: (\textbf{a}) 2016 RGB in \hl{5-6-4} Bands, (\textbf{b}) 2016 RGB in \hl{7-5-4} Bands, (\textbf{c}) 2016 unsupervised classification, (\textbf{d}) 2019 RGB in \hl{7-6-4} Bands, (\textbf{e}) 2019 RGB in \hl{7-5-3} Bands, (\textbf{f}) 2019 RGB: Clustering.}\label{fig08}
\end{figure}

Figures \ref{fig07}, \ref{fig08}, and \ref{fig09} show the wadi streams with their typical dendrite structure with comparison of images for \hl{the years }2013, 2015, 2016, 2019, 2022 and 2023. The fluvial network structure is oriented towards the Qena Bend segment of the Nile. The system of Wadi Fluvial structure and drainage stream network in the surroundings of the Qena Bend of the Nile River was identified on the images due to \hl{its }characteristic morphological structure \hl{with a }typical dendritic pattern that represent the flow direction of the dry wadi streams. Those are largely controlled by the underlying geological structure and regional lithology. The origin of these basins takes place in \hl{the }mountainous highland of the basement and is oriented West-East, either to the Nile or to the Red Sea, according to catchment, \hl{which means that }their direction is following these orientations. In both cases, the potential damage from  disastrous events may affect urban areas of the small cities and \hl{the infrastructure of industrial locations}. 

\begin{table}[H]
\small
\caption{Pairwise comparison of the k-means classification for Bands 5-4-3 of the Landsat-8 OLI/TIRS images for 2016 and 2019 in Qena Bend region. \textsuperscript{1}.\label{tab4}}
	\begin{adjustwidth}{-\extralength}{0cm}
		\begin{tabularx}{\fulllength}{||C?C|C|C?C|C|C||}
			\toprule
			\multirow{2}{*}{$Class ID$} & \multicolumn{3}{c}{\textbf{2016}} & \multicolumn{3}{c}{\textbf{2019}} \\
			\textbf{} & \emph{Band 5} & \emph{Band 4} & \emph{Band 3} & \emph{Band 5} & \emph{Band 4} & \emph{Band 3} \\
			\midrule
			1 & 25067.53 & 21805.82 & 18283.65 & 20721.42 & 18438.42 & 15725.33 \\
			2 & 19837.27 & 17506.45 & 15060.45 & 24198.84 & 21005.26 & 17763.84 \\
			3 & 27148.44 & 23483.89 & 19596.67 & 25536.93 & 22109.73 & 18481.20 \\
			4 & 16049.43 & 13237.00 & 11880.43 & 18257.00 & 9376.75 & 9707.50 \\
			5 & 19541.25 & 9323.50 & 9742.75 & 14476.67 & 12599.83 & 11411.17 \\
			6 & 12780.50 & 11767.00 & 10670.50 & 16992.00 & 15236.88 & 13304.75 \\
			7 & 21370.47 & 18891.65 & 16208.94 & 18829.53 & 16928.47 & 14654.00 \\
			8 & 23122.89 & 20069.28 & 16927.33 & 31113.00 & 26355.00 & 21077.00 \\
			9 & 18240.88 & 16536.62 & 14205.25 & 27920.00 & 23667.00 & 19115.00 \\
			10 & 17269.14 & 15186.43 & 13284.43 & 22736.25 & 19858.25 & 16938.06 \\	
			\bottomrule
		\end{tabularx}
	\end{adjustwidth}
	\noindent{\footnotesize{\textsuperscript{1} Clusters are computed for each band in the composite triplets: Band 5 (NIR), Band 4 (Red) and Band 3 (Green)}}
\end{table}

Figure \ref{fig09} visualises the Band combinations \hl{6-3-1, 5-4-3, 5-6-2 and 7-5-3} for the images covering Qena region in 2022 and 2023 years. 

\begin{figure}[H]
\makebox[1.00\linewidth][r]{%
	\begin{subfigure}[b]{.40\textwidth}
		\centering
			\includegraphics[width=0.87\textwidth]{Fig_09a}
		\caption{2022: Bands \hl{6-3-1}}
	\end{subfigure}%
	\begin{subfigure}[b]{.40\textwidth}
		\centering
		\includegraphics[width=0.87\textwidth]{Fig_09b}
		\caption{2022: Bands \hl{5-4-3}}
	\end{subfigure}%
	\begin{subfigure}[b]{.40\textwidth}
		\centering
		\includegraphics[width=0.90\textwidth]{Fig_09c}
		\caption{2022: Clustering}
	\end{subfigure}%
}\\
\vfill \vspace{1mm}
\makebox[1.00\linewidth][r]{%
	\begin{subfigure}[b]{.40\textwidth}
		\centering
			\includegraphics[width=0.87\textwidth]{Fig_09d}
		\caption{2023: Bands \hl{5-6-2}}
	\end{subfigure}%
	\begin{subfigure}[b]{.40\textwidth}
		\centering
		\includegraphics[width=0.87\textwidth]{Fig_09e}
		\caption{2023: Bands \hl{7-5-3}}
	\end{subfigure}%
	\begin{subfigure}[b]{.40\textwidth}
		\centering
		\includegraphics[width=0.92\textwidth]{Fig_09f}
		\caption{2023: Clustering}
	\end{subfigure}%
}\caption{Qena Bend of the Nile River: false colour composites and clustering of the Landsat 8-9 OLI/TIRS images: (\textbf{a}) 2022 RGB in \hl{6-3-1} Bands, (\textbf{b}) 2022 RGB in \hl{5-4-3} Bands, (\textbf{c}) 2022 unsupervised classification, (\textbf{d}) 2023 RGB in \hl{5-6-2} Bands, (\textbf{e}) 2023 RGB in \hl{7-5-3} Bands, (\textbf{f}) 2023 RGB: Clustering.}\label{fig09}
\end{figure}

The composite of \hl{5-4-3 }includes the NIR, Red and Green Bands that best represents the vegetation coverage in the Nile River Floodplain which is identified as dark crimson, well \hl{contrasted }against the other land cover types. 

\begin{table}[H]
\small
\caption{Pairwise comparison of the k-means classification for Bands 5-4-3 of the Landsat-8 OLI/TIRS images for 2022 and 2023 in Qena Bend region. \textsuperscript{1}.\label{tab5}}
	\begin{adjustwidth}{-\extralength}{0cm}
		\begin{tabularx}{\fulllength}{||C?C|C|C?C|C|C||}
			\toprule
			\multirow{2}{*}{$Class ID$} & \multicolumn{3}{c}{\textbf{2022}} & \multicolumn{3}{c}{\textbf{2023}} \\ 
			\textbf{} & \emph{Band 5} & \emph{Band 4} & \emph{Band 3} & \emph{Band 5} & \emph{Band 4} & \emph{Band 3} \\
			\midrule
			\midrule
			1 & 24428.50 & 21422.20 & 18051.50 & 21513.50 & 10245.750 & 10483.000 \\
			2 & 19238.75 & 17313.00 & 14911.25 & 25528.40 & 22234.800 & 18665.200 \\
			3 & 27402.67 & 23631.83 & 19472.50 & 17918.11 & 9096.444 & 9575.333 \\
			4 & 22303.50 & 19676.83 & 16793.33 & 24549.75 & 21362.375 & 17997.875 \\
			5 & 17517.33 & 11594.67 & 10989.33 & 15965.40 & 14835.200 & 13574.400 \\
			6 & 12943.67 & 11050.00 & 10303.00 & 22865.67 & 19456.933 & 16537.733 \\
			7 & 17225.22 & 15122.22 & 13253.00 & 23488.23 & 20695.000 & 17726.000 \\
			8 & 25640.82 & 22323.82 & 18730.73 & 18460.75 & 16531.833 & 14379.417 \\
			9 & 20699.44 & 18365.81 & 15840.12 & 20155.57 & 18063.786 & 15478.143 \\
			10 & 23487.22 & 20428.72 & 17181.22 & 27425.50 & 23626.900 & 19359.200 \\	
			\bottomrule
		\end{tabularx}
	\end{adjustwidth}
	\noindent{\footnotesize{\textsuperscript{1} Clusters are computed for each band in the composite triplets: Band 5 (NIR), Band 4 (Red) and Band 3 (Green)}}
\end{table}

The composite 5-6-2 includes the SWIR 1 and Blue Band which results in sands of the Western and Eastern Deserts looking in green shadows. This enables to discriminate them from the wadi channels that are coloured as light lilac colour. The limestone foothills are visualised by light magenta, while gritty and gravel soil -- by slate blue. Sandy soils are coloured by light chartreuse green, while the terrace soils and stony debris -- by lime green.

%%%%%%%%%%%%%%%%%%%%%%%%%%%%%%%%%%%%%%%%%%
\section{Discussion}

The Qena Wadi region is affected by annual flash floods and related geomorphological processes such as erosion as a consequence of the soil mass movement and deposition of soil particles during the heavy rainfalls. In turn, the deterioration of soil and changes in land cover types \hl{have a negative impact on }agriculture and may damage infrastructure, e.g., \hl{cut off roads. The consequences of such natural hazards emphasize }the dependence of social sustainability and urban infrastructure on the environmental setting of the region. Accurate information regarding the possibility of geomorphic risks or landslide deposition may be derived from \hl{an }evaluation of the landscapes's morphology. For instance, the evaluation of the channel incision modelled by geomorphic adjustment scenarios \hl{can demonstrate } the effects from geomorphic processes on infrastructure \cite{Hermas}.

\hl{Our study aimed to visualise gradual land cover changes in Upper Egypt and illustrate the geomorphology and topography of the Wadi Qena region. To achieve this, we employed a combination of methods including classification of }time series of \hl{satellite images, extraction of 2D and 3D mapping information from } topographic maps (GEBCO/SRTM grids)\hl{, }and \hl{evaluation of changes in land cover types using }statistical data derived from \hl{image analysis. The integration of such data as well as the analysis of the geological, geomorphological and climate setting based on prior works was used to evaluate the variations in land cover types illustrating the risks of hazards in the Qena Bend area. To further support our analysis, we integrated } geological, geomorphological \hl{, and climate data from prior works. By combining and analysing this data, we were able to evaluate }variations in land cover types \hl{and illustrate the hazards risks }in the Qena Bend area. \hl{Similar examples of data integration for hydrogeological modelling are presents in earlier studies focused on evaluating the aquifers potentiality in Eastern Desert} \cite{Moneim,ELHAMEED2017218}, 2D and 3D modelling for land surface slope gradient assessment \cite{Othman}, recognising the structural lineaments and fluvial channels and catchment analysis \cite{rs12101559}. \hl{These studies all demonstrate approaches to data fusion that benefit from advancements in computer vision algorithms for extracting} information from remote sensing data \cite{AZZAZY2022104805}.

Computer vision algorithms \hl{for }image recognition, discrimination\hl{, }and data processing are fundamental \hl{problems }in remote sensing and Earth sciences\hl{, with }a wide range of applications 
\cite{ABDELKAREEM2015245,Lemenkova202231,ELQADI2020104696,Husseinetal1995,DHANYA2022211}. The use of computer vision algorithms for the visualization, analysis, and classification of satellite images \hl{aims to prevent }the risk of natural hazards \hl{by mapping }potentially endangered areas. \hl{Regions such }as wadi valleys in desert areas of southern Egypt are at high \hl{risk }of flash floods. \hl{In this paper, we present }an approach to represent, process, \hl{analyze, }and compare satellite images, both as classified maps and as extracted information on land cover types and their changes as statistical data derived from the images and \hl{summarized}in tables. We \hl{demonstrate }how an R-based approach\hl{, utilizing }k-means clustering and spectral reflectance characteristics of the remotely sensed data\hl{, }can be used to learn how landscapes in \hl{the }southern Egypt gradually change over time.

\hl{To recognize changes in the satellite images, }short-time series analysis \hl{was performed }using principles of remote sensing data processing \cite{curran1985principles}. \hl{Similarity }in texture and structure of the objects represented on the Landsat 8-9 scenes was \hl{analysed }based on the characteristics of the multispectral channels from the Landsat satellite images using fusion techniques for \hl{diverse }spectral bands \cite{4234117}. \hl{False }colour composites (NIR, Red\hl{, }and Green) \hl{were }applied to better \hl{recognize } the attributes of vegetation, water\hl{, }and land as major classes by automatic machine-based algorithms of R. \hl{As a result, } subclasses such as wadi valleys, Nile floodplain and valley, \hl{and hills }were defined according to the surface structure in the study area. \hl{The } orientation of the patterns of Wadi Qena \hl{on the satellite images is }almost parallel to the axis direction of the Red Sea uplift, which corresponds with the \hl{geological} origin of the wadi network.

\hl{A pairwise comparison of the land cover changes enables to make an  general picture of the gradual trend in land cover changes in Qena surroundings over the 10 years. Figures} \ref{fig07}, \ref{fig08} and \ref{fig09} \hl{present a pairwise comparison of the changes in land cover types in Qena Bend region up to the relevant two years next to each other. The analysis of these images demonstrate maps of changes in land cover types calculated and visualised between the years 2013 and 2015 (Figure} \ref{fig07}) \hl{and the years 2016 and 2019 (Figure} \ref{fig08}) \hl{and the years 2022 and 2023 (Figure} \ref{fig09}). \hl{The computation is summarised in the tables in Table} \ref{fig03} \hl{for 2013/2015, Table} \ref{fig04} \hl{for 2016/2019, and Table} \ref{fig05} \hl{for 2022/2023.} 

Flash floods and soil deposition in \hl{the }Wadi Qena region are \hl{natural phenomena related to the climate and occur repeatedly in } southern Egypt due to specific geomorphology \cite{ElShamy}. Therefore, predicting such hazards \hl{requires }complex climatic computational modelling. \hl{However, }robust information derived from remote sensing data processed by advanced computer vision and pattern recognition \hl{tools can support risk mitigation measures }and propose planning strategies in regions prone to natural hazards.

\hl{Geomorphic data can also be used to evaluate groundwater supply }\cite{Abdelkhalek2021,Abdelkareem} and its availability for irrigation in agriculture \cite{Kamal2021}.
Therefore, the satellite-derived \hl{maps of }land cover types are a useful tools for \hl{both }determining hazard areas and defining zones susceptible to risks of floods\hl{, as well as }applied geomorphic modelling according to the distribution of fluvial network of wadi in \hl{the } southern segment of the Eastern Desert. Moreover, data analysis and visualization performed using scripting methods provide a new source of information that can be used for taking mitigation policies and measures of sustainable development in \hl{the }Qena Band region, as well as urban planning in hazard prone areas of Egypt \cite{ABDELAAL2020103671}. 

\hl{In }future works, \hl{we plan to assess the }geomorphological hazards by \hl{analysing the }susceptibility of various regions supported by the classified satellite images\hl{. This will enable the evaluation of the risk of flash floods} , soil erosion\hl{, and mass deposition. Since } landscapes are subject to destructive geomorphological hazards such as flash floods, dust storms\hl{, }and heavy showers, \hl{the use of satellite images provides }objective and robust information regarding these hazards and \hl{their }disastrous processes. \hl{Flash floods, in particular, cause sediment deposition }and debris and affect \hl{the }shape of the land surface\hl{. Therefore, }preventive mapping and evaluation of \hl{geomorphological risks are crucial }for infrastructure and agricultural activities in Upper Egypt.

%%%%%%%%%%%%%%%%%%%%%%%%%%%%%%%%%%%%%%%%%%
\section{Conclusions}

\hl{The use of satellite images in the classification of land cover types by clustering is an important approach for applications in Earth sciences. At the same time, identification of the features and properties of the landscapes on the space born images is a challenging field of investigation in geoinformation due to the complexity of the interpretation of the land cover classess. A multi-disciplinary approach combining the remote sensing data and advanced technical tools of their processing is an effective approach that supports satellite image analysis. In this paper, we demonstrated the application of such an approach by presenting the programming tools used for spatial data analysis. We showed that using the advanced methods of scripting languages for image processing is an effective tools for image analysis and classification.}

\hl{Specifically}, we have demonstrated a method \hl{for extracting }information from the Landsat 8-9 OLI/TIRS \hl{satellite images }to track land cover changes in wadi drainage surface of \hl{the }Qena Bend district of the Nile River. We presented an application of R scripting libraries for processing geoinformation to analyse visual objects on the remote sensing data by computer vision and methods of pattern recognition on the images. Local features in the target study area were identified using k-means clustering method and visualized \hl{in synthesised }maps based on the \hl{six classified }images. This algorithm exploits the values of the pixels' DNs to discriminate the classes on the land surface based on the similarity of pixels' values and assigning them to 10 \hl{separate }classes. 

We have approached the problem of unsupervised classification of the satellite images \hl{through the use of the R language, for extracting objects from }information on the land cover types \hl{in the }Qena Bend region \hl{in }Upper Egypt. \hl{The R scripting }libraries 'raster', 'terra'\hl{, }and 'RStoolbox' \hl{were utilised to extract and analyse information from }remote sensing data \hl{in the context of the geomorphological structures over the studied area. Through a machine-based approach enabled by R, spatial features and land cover classes were automatically separated, allowing for the analysis of changes in the ephemeral riverbed complex in the Qena district using }computer vision and applied programming.

In the framework of this study, the \hl{scripted }implementation of the R libraries \hl{was }evaluated under Landsat OLI/TIRS data with auxiliary cartographic data processing by GMT. The comparison of the land cover changes \hl{was then summarised in tabular form, based on }information extracted from the R-processed images. Our analyses shows that \hl{the R approach }algorithms of computer vision applied for geographic data processing is effective for multispectral \hl{space-borne data. Using R, we have computed maps for pairwise comparison of the different target years (2013 and 2015, 2016 and 2019, 2022 and 2023) to obtain a general trend of the transformation in selected land cover types in Qena region.}

Our study demonstrates the effectiveness of the R-based method \hl{as }an advanced approach to geospatial image processing compared to state of the art GIS methods. \hl{Specifically, our algorithm is }trained to classify pixels and assign them to \hl{different }land cover classes \hl{based on the spectral reflectance values }in various bands of the satellite images\hl{, allowing for automatic interpretation of the images. The use of }computer vision algorithms \hl{and automation allowed for the generation of } a series of \hl{maps depicting }the land cover types in \hl{the Qena area during the years }2013, 2015, 2016, 2019, 2022\hl{, }and 2023. \hl{Comparing the }changes in land cover types\hl{provides a better understanding and interpretation of }the cumulative effects \hl{of } environmental changes and occasional floods in \hl{the }southern segment of the Eastern Desert on the geomorphic patterns \hl{of } land cover types \hl{in the Qena Bend region of Upper Egypt}.

%%%%%%%%%%%%%%%%%%%%%%%%%%%%%%%%%%%%%%%%%%
\authorcontributions{Supervision, conceptualization, methodology, software, resources, funding acquisition, and project administration, O.D.; writing---original draft preparation, methodology, software, data curation, visualization, formal analysis, validation, writing---review and editing, and investigation, P.L. All authors have read and agreed to the published version of the manuscript.}

\funding{The publication was funded by the Editorial Office of \emph{Information}, MDPI, by providing 100\% discount for the APC of this manuscript. This project was supported by the Federal Public Planning Service Science Policy or Belgian Federal Science Policy Office - BELSPO (B2/202/P2/SEISMOSTORM).}

\institutionalreview{Not applicable.}

\informedconsent{Not applicable.}

\dataavailability{Not applicable} 

\acknowledgments{The authors thank the reviewers for reading and review of this manuscript.}

\conflictsofinterest{The authors declare no conflict of interest.} 

\abbreviations{Abbreviations}{
The following abbreviations are used in this manuscript:\\

\noindent 
\begin{tabular}{@{}ll}
ASTER & Advanced Spaceborne Thermal Emission and Reflection Radiometer \\ 
DBSCAN & Density-Based Spatial Clustering of Applications with Noise \\
GEBCO & General Bathymetric Chart of the Oceans \\
GMT & Generic Mapping Tools \\
GIS & Geographic Information System \\
Landsat OLI/TIRS & Landsat Operational Land Imager Thermal Infrared Sensor \\
NIR & Near Infrared \\
RGB & Red Green Blue\\
SRTM & Shuttle Radar Topography Mission \\
SWIR & Shortwave Infrared \\
DEM & Digital Elevation Model\\
UTM & Universal Transverse Mercator \\
WGS84 & World Geodetic System 1984 \\
\end{tabular}
}

\begin{adjustwidth}{-\extralength}{0cm}
%\printendnotes[custom] % Un-comment to print a list of endnotes

\bibliography{Qena.bib}
\reftitle{References}
\nocite{*}

%=====================================
% References, variant A: external bibliography
%=====================================


%%%%%%%%%%%%%%%%%%%%%%%%%%%%%%%%%%%%%%%%%%
\end{adjustwidth}
\end{document}
